\hypertarget{group-state-and-general-measurement}{%
\section{Group state and general measurement}\label{method-insurance-methods-group-state-and-general-measurement}}

The following cases apply selected general-measurement and service-allocation provisions from the retained AASB 17 compilation \citep{aasb17}. They extend the premium-allocation example without presuming actuarial calibration or a complete insurer model.

The general method distinguishes fulfilment cash flows (FCF), including an
explicit risk adjustment (RA), from the contractual service margin (CSM),
which holds unearned profit. The liability for remaining coverage (LRC)
concerns future service; the liability for incurred claims (LIC) concerns
service and insured events already experienced. Keeping these components
separate makes a change in future estimates different from an unpaid claim.

Contract records retain insured risks, premiums/benefits, service and settlement dates, options/guarantees, cancellation/repricing rights and acquisition costs. Group admission evaluates portfolio/similar-risk management, onerous status and required issuance cohorts under AASB 17.14--24, recognition under 25--28 and boundary under 34. Cash outside the boundary belongs to future contracts, not this group's profit. GroupState retains expected dated cash flows, current discount curves, inception rates, explicit risk adjustment, CSM, loss component, LRC/LIC, coverage units, service experience and assumption vintages. Same premium totals are not sufficient state for merging different guarantees or claim tails.

General measurement follows 32--38: fulfilment cash flows are probability-weighted future net outflows, adjusted for timing/financial risk without double counting, plus explicit non-financial risk adjustment. CSM removes initial unearned gain under its conditions; an initial net loss is recognized under 47, not negative CSM. Paragraphs 40--44 separate remaining coverage and incurred claims. For a nonparticipating group, accrete CSM at the specified locked-in rate, apply future service changes at their required rate to CSM up to its available amount, record any excess loss and allocate current service over current plus future coverage units under B119. B96--B97 distinguish future service from incurred-claim changes, financial changes and current experience; not every revision is a CSM adjustment.

A MeasurementInput must resolve cash-flow/law, discount, risk-adjustment method, service grouping and coverage-unit evidence. Missing probabilities or RA support permit conditional sensitivity cases but block a source-supported measurement. Risk adjustment is the compensation the entity requires under 37; a standard deviation, confidence quantile or arbitrary percentage is only a candidate method until justified for that entity. Disclosures and the confidence-level equivalent, where applicable, are separate obligations. Market-consistent financial inputs and entity risk compensation must not be conflated.

\hypertarget{complete-two-year-general-measurement-case}{%
\section{Complete two-year general-measurement case}\label{method-insurance-methods-complete-two-year-general-measurement-case}}

Take a two-year nonparticipating group with one premium 120 received at inception, claims expected/actually incurred and paid 40 each year, zero discount rate, RA 10 initially released 5 with each equal service year, no acquisition costs, investment component, lapses, taxes, finance changes or reinsurance. Coverage units are 1 in each year because equal benefits and coverage are assumed. Before premium payment, expected net cash outflow is 80\ensuremath{-}120=\ensuremath{-}40, RA 10, FCF \ensuremath{-}30 and CSM 30, with net initial carrying amount zero. After the premium debit cash 120/ credit LRC 120. Its components are remaining expected claims 80, RA 10 and CSM 30.

Each year releases expected service claims 40, RA 5 and CSM 15: debit LRC 60/ credit insurance revenue 60. Incur claims 40: debit service expense 40/credit LIC 40; payment debits LIC 40/credits cash 40. End of year one: cash 80, LRC 60 (claims 40, RA 5, CSM 15), LIC 0 and equity/profit 20. End year two: cash 40, liabilities 0, cumulative equity/profit 40. Revenue 120 and claims 80 reconcile to premium 120 less claims cash 80. No investment component has been mistaken for revenue.

At year-one close, a separate pre-next-service change increases the remaining claim estimate from 40 to 60, RA unchanged. FCF rises 20; consume remaining CSM 15 and recognize service loss 5. LRC rises from 60 to 65, comprising FCF 65, CSM 0 and loss component 5 as a tracked portion, not an added liability on top. Debit service loss 5/credit LRC 5 with component transfer CSM\ensuremath{\rightarrow}FCF 15. A decrease of 8 in the revised future claim estimate, before any additional service/finance movement, first reverses loss 5 and restores CSM 3. FCF 57 plus CSM 3 gives LRC 60. Credit loss reversal 5/debit LRC 5. This tests 48--52 ordering. Later systematic allocation of loss component to releases/finance must reduce it to zero by coverage end; it cannot be amortized as revenue a second time.

In a separate case, an already-incurred unpaid claim estimate rises 40 to 48. Debit service expense 8/credit LIC 8; do not consume future-service CSM. A discount-rate change is financial under B97; its profit/OCI route follows the selected permitted finance policy. The zero-rate base case does not remove the need for separate current and locked-in rate calculations in a real group.

For the unreversed onerous alternative, carry the loss component through the remaining service year. Opening LRC 65 comprises future claims 60 and RA 5, with CSM 0 and loss component 5 within that LRC. A declared proportional allocation assigns 5/65 of each expected claim/RA release to the loss component, consistent with paragraphs 49--52 and B124's revenue exclusions. The expected claim release allocates 60/13 to loss reversal and the RA release 5/13; the total is 5, leaving loss component 0. Debit LRC 65, credit revenue 60 and credit loss reversal 5. Separately incur and pay claims 60 through LIC. Period profit is 60\ensuremath{-}60+5=5; combined with first-year profit 20 and the earlier onerous loss 5, lifetime profit is 20. Closing cash 20 equals equity 20 and all liabilities are zero. Treating all 65 as revenue while also reversing loss 5 overstates profit by 5.

\hypertarget{nonzero-discount-and-service-allocation}{%
\section{Nonzero discount and service allocation}\label{method-insurance-methods-nonzero-discount-and-service-allocation}}

A second complete case keeps premium 120, two equal coverage units and RA 10, but expected and realized claims are 44 after one year and 48.4 after two years. Both the inception and current discount curves remain a flat 10\% throughout; there are no new estimates, options, fees, taxes, investments or OCI election. This RA schedule is an authored measurement input, not an actuarial calibration: it declines by 5 per year, with its entire change in the service result under 81. All other insurance finance effects go to profit or loss. The retained source supports the mechanics in 44, 80--88,B119 and B123--B124.

At inception the present value of the two expected claims is
\[
 \frac{44}{1.1}+\frac{48.4}{1.1^2}=40+40=80.
\]
 FCF before premium is 80+10\ensuremath{-}120=\ensuremath{-}30, CSM 30, and after premium LRC 120. In year one the claim PV accretes 8; CSM accretes 3 at the locked-in rate. Allocate the resulting CSM 33 across the two current/future units, releasing 16.5. Revenue is expected claims 44 plus RA 5 plus CSM 16.5=65.5. Service expense 44 and finance expense 11 leave profit 10.5. After paying 44, cash 76 and LRC 65.5 (claim PV 44, RA 5, CSM 16.5) reconcile to equity 10.5.

In year two claim PV accretes 4.4 and CSM accretes 1.65. Release all remaining CSM 18.15, RA 5 and expected claims 48.4, giving revenue 71.55. Service expense 48.4 and finance expense 6.05 leave profit 17.1. Final cash/equity 27.6 equals premium 120 minus claims 92.4; all group liabilities are zero. Total revenue 137.05 less finance expense 17.05 equals premium 120. Setting revenue equal to cash premium each year or releasing 15 of CSM each year would ignore the specified financing. These are separate component calculations followed by debit finance expense/ credit LRC for accretion, debit LRC/credit revenue for service and the stated LIC recognition/settlement entries. Cash is deliberately assumed uninvested.

\hypertarget{paa-reinsurance-and-valuation-methods}{%
\section{PAA, reinsurance and valuation methods}\label{method-insurance-methods-paa-reinsurance-and-valuation-methods}}

The preceding premium-allocation case states separate eligibility and recognition timing for gross insurance and reinsurance. PAA under 53--59 requires the one-year coverage or nonmaterial-difference condition; a short payment lag is not the eligibility test. LRC starts with received premium and applicable acquisition adjustments, then releases insurance revenue under the prescribed service pattern. LIC still measures incurred obligations using the applicable fulfilment-cash-flow rules; the allowed short-settlement discount simplification needs its own conditions. Emerging onerous facts require a loss assessment; a premium balance cannot be presumed sufficient.

Reinsurance has its own group, boundary, recognition time, counterparty risk, cash-flow, CSM/cost and possibly loss-recovery component under 60--70A. Do not multiply every gross insurance liability by a quota share: excluded claims, limits, deductibles, timing, reinstatements and counterparty non-performance matter. Claim 100 with an eligible 50\% recovery but expected 90\% collectability gives conditional expected recovery 45, not 50; altering transferred insurance risk and reinsurer non-performance are different inputs. This simple calculation does not establish treaty recognition or RA by itself. Gross liability 100 remains; reinsurance recovery is a separate asset/measurement with its own authority.

For direct participation, test B101's conditions using the actual underlying items, substantial-share and variation judgments. Retain underlying fair values, entity fee share and guarantee/risk-mitigation choices; .45 and B101--B118 specify the variable-fee adjustments to CSM, with onerous and mitigation exceptions. A local component illustration of a qualifying share: underlying-value increase 20 and entity's fixed 10\% share gives candidate variable-fee change 2; eligibility, guarantees and other FCF changes must be evaluated before posting. The full option/guarantee valuation and transition calculations are further supported profiles, with source and numerical obligations, not an assumed linear fee model.

The claims-development example above supplies an estimation input, not a
recognition rule. Its point projection is 148, but the single observed final
factor leaves its variance unresolved. In the general method the factors are
\[
 \widehat f_k=\frac{\sum_{i\in O_k} C_{i,k+1}}{\sum_{i\in O_k}C_{i,k}},
\]
where $O_k$ contains origin years with both required observed cells. Each
incomplete row is projected through the remaining factors. The same triangle
therefore has ultimate claims 180, 216 and 252, with unpaid amounts 0, 36 and
112. Tail development, calendar inflation, zero or negative cells and dependence
need explicit extensions. The conditional means, variance scaling and
independent-origin-year assumptions behind the inspected Mack error formulas
\citep[equations (1)--(3), Theorems 1--3]{mack1993} cannot turn this tiny
triangle into a predictive law or a justified insurance risk adjustment.

\hypertarget{interfaces-exclusions-and-acceptance-cases}{%
\section{Interfaces, exclusions and acceptance cases}\label{method-insurance-methods-interfaces-exclusions-and-acceptance-cases}}

Insurance events are group\_recognize, premium\_settle, service\_release, claim\_incur, claim\_settle, assumption\_change, onerous\_change and treaty\_change, each bound through the event-basis, authorization and commit contracts to the exact group/version, evidence and information view. Journals and component roll-forwards must reconcile independently; a component reallocation is not always a net liability change. Invalid grouping, missing rate/law, unsupported option valuation, ambiguous contract boundary or unknown service units produces the corresponding unresolved/outside-scope result. Failed numerical paths retain their coverage/mass, without successful-only normalization. Modified contracts need the applicable derecognition analysis; no new ID can erase unfavorable old-group obligations.

Initial transition requires an actual method choice and eligibility for full retrospective, modified-retrospective or fair-value transition under Appendix C, with acquisition, grouping, past service and historical-rate records. This collection does not claim an executed transition or a general VFA valuation. The nonzero-rate general case and onerous allocation through expiry above are finite specification examples. Required later acceptance still includes a priced guarantee, acquisition/reinsurance interaction, transition and actuarial backtesting. These retained obligations prevent the two-year illustration from being described as full insurance-domain readiness.
